\documentclass[11pt]{article}
\usepackage[a4paper,margin=26mm]{geometry}
\usepackage{amsmath,amssymb,amsthm,mathtools,hyperref,microtype}
\hypersetup{colorlinks=true,urlcolor=blue,linkcolor=blue}
\newtheorem{proposition}{Proposition}
\newtheorem{lemma}[proposition]{Lemma}
\newcommand{\Ric}{\operatorname{Ric}}
\newcommand{\Rm}{\operatorname{Rm}}
\newcommand{\CP}{\mathbb{CP}}
\newcommand{\vol}{\operatorname{Vol}}
\title{Smooth attainment and a genuine bubbling sublevel\\on a Godeaux blowup}
\author{Research continuation}
\date{8 October 2026}
\begin{document}
\maketitle
\begin{abstract}
The existence target is a smooth global minimizer of
$P(g)=\int|\Ric_g^\circ|^2$ which is not Bach-flat. On the fixed classical
Godeaux blowup, the supplied strict Burns insertion theorem gives
$p_*<8\pi^2$. We exhibit an actual Eguchi--Hanson degeneration on this
same smooth manifold at a level strictly below $8\pi^2$. Thus this bound
is not a compactness threshold. We also give one aggregate characteristic
class estimate for marked negative-definite splittings, identifying precisely
why it does not detect this degeneration. These conclusions do not establish
attainment or nonattainment of $p_*$.
\end{abstract}

\section{The existence question}
Let $X$ be the quotient of the smooth Fermat quintic by the free action
\[
 \gamma[x_1:x_2:x_3:x_4]
 =[\zeta x_1:\zeta^2x_2:\zeta^3x_3:\zeta^4x_4],\qquad \zeta^5=1,
\]
where $\zeta$ is primitive, and put $M=X\#\overline{\CP}^{\,2}$.
Then $K_X^2=1$, $\chi(M)=12$, $\tau(M)=-8$, and $T=2\chi+3\tau=0$.
The inherited classical monopole estimates and the project's strict
negative K\"ahler--Einstein insertion theorem give
\begin{equation}\label{eq:interval}
 \frac{16\pi^2}{3}\le p_*(M):=\inf_{g\in\operatorname{Met}(M)}P(g)<8\pi^2.
\end{equation}
Here the strict upper bound is an input from the supplied manuscript,
not independently re-proved in this note. The exact input is its general
KE expansion, at one insertion point,
\[
 \frac{P(g_\varepsilon)}{8\pi^2}
 =1-\frac{s^2}{64}\varepsilon^2\log(1/\varepsilon)+O(\varepsilon^2).
\]
It is supported in a regular coordinate ball and keeps the outside metric
unchanged; see the source record at the end.

A smooth minimizer would have constant scalar curvature by conformal
variation. All conformal classes here have negative Yamabe constant by
the actual monopole input. A Bach-flat $P$-critical metric of nonzero
constant scalar curvature is Einstein, by the two Euler--Lagrange equations.
But Einstein and $T=0$ imply $s=W^+=0$ by Chern--Weil, contradicting the
scalar estimate $\int s^2\ge32\pi^2$. Thus smooth attainment on this $M$
would solve the desired non-Bach-flat minimizer problem.

The inference missing from \eqref{eq:interval} is smooth attainment.
Anderson's ideal-minimizer theorem \cite{Anderson} admits orbifold and
open/collapsing alternatives; it does not remove them from
\eqref{eq:interval}.

\section{An explicit nodal Godeaux degeneration}
Define the invariant quintic
\begin{align*}
 F={}&2(x_1^5+x_2^5+x_3^5+x_4^5)\\
 &-x_1^3x_3x_4-x_1^2x_2^2x_4-x_1^2x_2x_3^2-x_1x_2^3x_3\\
 &-x_1x_2x_4^3-x_1x_3^2x_4^2-x_2^2x_3x_4^2-x_2x_3^3x_4.
\end{align*}
Every monomial has weight zero modulo five. At $[1:1:1:1]$ the gradient
vanishes and the homogeneous Hessian is
\[
 40I_4-10\mathbf 1\mathbf 1^{\mathsf T}.
\]
Its only kernel is the radial line, so the projective singularity is a node.
For completeness, exact polynomial reduction gives, in the chart $x_4=1$,
\[
 \left(\partial_{x_i}F\big|_{x_4=1}:1\le i\le4\right)
 = (x_1-x_3^3,\ x_2-x_3^2,\ x_3^5-1)
 \quad\hbox{in }\mathbb Q[x_1,x_2,x_3].
\]
On $x_4=0$, in each of the three charts $x_i=1$, the derivative ideal is
the unit ideal. Hence the hypersurface has precisely five nodes, a single
free $\mathbb Z_5$ orbit. The derivative-ideal calculation and the Hessian
were verified by exact rational arithmetic in the accompanying script.

The only fixed points of nonidentity group elements in projective space
are the coordinate vertices; none lies on $F=0$. Therefore
\[
 Z_0=\{F=0\}/\mathbb Z_5
\]
has exactly one $A_1$ singularity. Its canonical bundle is ample:
the quotient pullback is $\mathcal O(1)$ by adjunction.

The invariant pencil $F+u\sum x_i^5$ smooths the nodes for generic small
$u\ne0$. The smooth invariant quintics avoiding the fixed vertices form
a connected complex Zariski-open set. Equivariant Ehresmann triviality
identifies their quotients with the same oriented smooth $X$.
The $A_1$ Milnor fibre is diffeomorphic to the disk bundle of degree $-2$;
equivalently, its smoothing and minimal resolution have the same oriented
smooth filling. Consequently $X$ contains a smooth sphere $C$ with
\[
 C^2=-2,\qquad K_X\cdot C=0.
\]
After a blowup outside this region, $E\cdot C=0$ as well. This is an
embedded sphere, rather than just a root of the intersection lattice.

\section{Bubbling strictly below the known upper threshold}
\begin{proposition}\label{prop:bubble}
Assume the supplied local strict Burns insertion theorem stated in
Section~1. There are smooth metrics $g_\lambda$ on the fixed manifold $M$
such that
\[
 P(g_\lambda)\longrightarrow q<8\pi^2,
\]
and a rescaled region converges to the nonflat Ricci-flat
Eguchi--Hanson metric. After volume normalization this sequence has no
smooth convergent subsequence modulo diffeomorphisms.
\end{proposition}
\begin{proof}
The ample canonical orbifold $Z_0$ admits a negative KE orbifold metric
$h_0$ by the orbifold Aubin--Yau theorem; see \cite{Kobayashi,Suvaina}.
Apply the supplied insertion at a regular point away from the node.
Fix its parameter sufficiently small to obtain an orbifold metric $h$
on the one-point blowup of $Z_0$ with $P(h)=q<8\pi^2$. It equals $h_0$
near the node.

Replace the node with Eguchi--Hanson at length scale $\lambda$. In shared
end coordinates the differentiated expansions are
\[
 h=\delta+O(r^2),\qquad
 g_{\mathrm{EH},\lambda}=\delta+O(\lambda^4r^{-4}).
\]
Interpolate at $\rho=\lambda^{1/2}$. The annulus has volume $O(\rho^4)$
and curvature bounded by $C(1+\lambda^4\rho^{-6})$, so
\[
 \int_{\mathrm{annulus}}|\Ric^\circ|^2
 \le C(\rho^4+\lambda^8\rho^{-8})\longrightarrow0.
\]
The removed Einstein orbifold ball and the inserted Ricci-flat interior
both have $P=0$. Thus $P(g_\lambda)\to q$. The resulting smooth manifold
is the smoothing of $Z_0$ blown up at a separate point, hence is $M$.
The total volume tends to the positive volume of $h$, whereas the maximum
Riemann curvature in the EH region grows as $\lambda^{-2}$. Uniform
volume normalization preserves this curvature divergence.
\end{proof}

This proposition makes no assertion that $q=p_*(M)$. It disproves the
suggested compactness of the whole sublevel below $8\pi^2$, not smooth
attainment itself. Its ALE interpolation is classical; the supplied strict
Burns saving is what puts the sequence below that threshold.

\section{What the actual classes control in an entire splitting}
The following estimate collects all marked components simultaneously.
It is conditional on an appropriate splitting and actual field passage;
it is not a claim that every minimizing sequence already supplies them.

\begin{proposition}\label{prop:aggregate}
Suppose a bulk $B$ and finitely many ALE pieces $N_i$ give a marked rational
orthogonal decomposition, with negative-definite core subspaces. Assume
relative Chern--Weil additivity $T_B+\sum_iT_i=0$ and actual bulk monopole
input for the restrictions of $K+E$ and $K-E$, sufficient for the classical
mixed curvature estimate. Write
\[
 K=K_B+\sum_iK_i,\qquad E=E_B+\sum_iE_i,\qquad
 d=-\sum_iK_i^2,\quad e=-\sum_iE_i^2.
\]
With the standard operator norm for Weyl curvature, put
\[
 C_i=\int_{N_i}\left(4|W^+|^2+\frac{s^2}{12}\right)\ge0.
\]
Then the total energy of the configuration satisfies
\begin{equation}\label{eq:aggregate}
 P_B+\sum_iP_i\ge
 \frac{16\pi^2}{3}\bigl(1+d+\max\{e-1,0\}\bigr)+\sum_iC_i.
\end{equation}
\end{proposition}
\begin{proof}
Since $K^2=1$ and $E^2=-1$, one has $K_B^2=1+d$ and $E_B^2=e-1$.
The average of the two self-dual projection squares is
\[
 \frac{|(K_B+E_B)^+|^2+|(K_B-E_B)^+|^2}{2}
 =|K_B^+|^2+|E_B^+|^2
 \ge1+d+\max\{e-1,0\}.
\]
Apply the classical inequality
$P_B\ge8\pi^2(\frac23\max_\pm|(K_B\pm E_B)^+|^2-T_B)$.
For each core, relative Chern--Weil gives $P_i=C_i-8\pi^2T_i$.
Addition cancels all the relative $T$ terms and yields
\eqref{eq:aggregate}.
\end{proof}

In particular, configurations of total energy less than $8\pi^2$ satisfy
\[
 d+\max\{e-1,0\}+\frac{3}{16\pi^2}\sum_iC_i<\frac12.
\]
For the actual $A_1$ sphere above, both restrictions vanish and the
Ricci-flat ASD core has $C_i=0$. The aggregate estimate thus retains only
$16\pi^2/3$. The missing information cannot be supplied by counting that
sphere again.

For comparison, $r$ disjoint degree-$-3$ cores with canonical degree one
give $d=r/3$, excluding $r\ge2$ at this energy. A chain
$[-3,-2,\ldots,-2]$ of rank $r$ with canonical degrees $(1,0,\ldots,0)$
instead has $d=r/(2r+1)<1/2$, from its inverse intersection matrix.
This last calculation does not assert that every such chain embeds in $M$.

\section{Attainment does not require compactness of every minimizing sequence}
There is also a concrete example within the construction above. On the
unblown smooth $X$, a negative KE metric exists and attains $\inf P=0$.
Resolving the node of $(Z_0,h_0)$ by shrinking EH metrics, without making
the Burns insertion, gives another sequence on $X$ with $P\to0$ and
unbounded curvature. Thus a smooth global minimizer and a degenerating
minimizing sequence coexist on the same smooth manifold.

Accordingly, the existence target for $M$ is to obtain at least one smooth
global minimizer. Strict comparison with all ideal boundary configurations
would be one sufficient route, if their compactification and regularity
were established; it is not a necessary extra theorem. A suitable selection
of minimizing sequences or a smooth minimizing construction could also
succeed. Proposition~\ref{prop:bubble} does not decide among these routes.

What is unresolved is a genuinely minimizing comparison: no result here
shows that a smooth competitor is lower than every relevant optimized
singular configuration, and no smooth metric is shown to attain $p_*(M)$.
The previous single-core mass calculations therefore remain tools for one
stratum, rather than a proof of Q3.

\paragraph{Provenance of the nonclassical input.}
The strict insertion expansion is taken from the supplied
\texttt{01\_Research\_Synthesis.tex}, subsection ``The global upper bound
that changed during the night'', and its referenced A02.59D source.
That subsection already extends the construction to canonical models
with ADE points. The explicit nodal quintic, its exact finite verification,
and the aggregate calculation above explain the consequence for this fixed
candidate. They are not assertions of a newly proved general attainment
theorem or an exhaustive originality claim.

\begin{thebibliography}{9}
\bibitem{Anderson} M. T. Anderson, \emph{Canonical metrics on
3-manifolds and 4-manifolds}, especially Theorem 4.10 and Remark 4.12,
\href{https://arxiv.org/abs/math/0511034}{arXiv:math/0511034}.
\bibitem{Kobayashi} R. Kobayashi, \emph{Einstein--K\"ahler V-metrics on
open Satake V-surfaces with isolated quotient singularities},
Math. Ann. \textbf{272} (1985), 385--398,
\href{https://doi.org/10.1007/BF01455566}{doi:10.1007/BF01455566}.
\bibitem{Suvaina} I. \c Suvaina, \emph{The Yamabe invariant of a class
of symplectic manifolds}, Proposition 2.4 and proof of Theorem 3.4,
\href{https://arxiv.org/abs/1410.1236}{arXiv:1410.1236}.
\bibitem{LeBrun} C. LeBrun, \emph{Ricci curvature, minimal volumes,
and Seiberg--Witten theory},
\href{https://arxiv.org/abs/math/0003068}{arXiv:math/0003068}.
\end{thebibliography}
\end{document}
